\chapter{Notación}\label{ap:notacion}

\begin{longtable}{lp{11.5cm}}
\toprule
Símbolo & Significado \\
\midrule
\endhead
$Z$, $i,j$ & Número de zonas (12) e índices de zona \\
$X_{i,t}$ & Índice de demanda motorizada de la zona $i$ en el año $t$ (MBG) \\
$\mu_i$, $\Sigma$ & Deriva y matriz de covarianza de los logaritmos del índice \\
$\sigma_c$, $s_i$ & Volatilidad del factor común y volatilidad idiosincrática de la zona $i$ \\
$\lambda$, $J$ & Intensidad y tamaño (lognormal) de los saltos de tipo Merton \\
$P_i$, $A_j$ & Viajes generados y atraídos en hora pico \\
$T_{ij}$ & Matriz origen--destino (modelo gravitacional doblemente acotado) \\
$\beta$ & Parámetro gravitacional \\
$c_{ij}$ & Costo compuesto (logsum) entre zonas \\
$p^{\text{TP}}_{ij}$ & Probabilidad logit de elegir transporte público \\
$x_e$ & Razón volumen/capacidad del corredor $e$ \\
$f(x)$ & Función de congestión BPR o de hacinamiento \\
$\theta$ & Intensidad de la política de proximidad y policentrismo \\
$d$ & Diseño urbano (conjunto de proyectos e intensidad $\theta$) \\
$C_t(d,\omega)$ & Costo social del diseño $d$ en el año $t$ y el escenario $\omega$ \\
$\VP(d,\omega)$ & Valor presente del costo social \\
$r$ & Tasa social de descuento (12\,\% en la fase I, 9\,\% en la fase II) \\
$\text{BN}(d)$ & Beneficio neto frente al diseño de referencia: $\VP(\text{ref})-\VP(d)$ \\
$\CVaR_\alpha$ & Valor en riesgo condicional al nivel $\alpha=0{,}95$ \\
$w$ & Peso del CVaR en el criterio media--CVaR \\
$E^*$ & Costo urbano anual del viaducto que iguala elevado y subterráneo \\
\bottomrule
\end{longtable}
